% 第6章 ボルツマン方程式
\section{ボルツマン方程式}
\label{ch:ch06_boltzmann}

\paragraph{この章で導くこと（入力→出力）}
一般のボルツマン方程式 $df/d\lambda=C[f]$（入力）から、(a) 光子の温度ゆらぎ階層、
(b) CDM 無衝突流体、(c) バリオン流体（運動量交換）、(d) ニュートリノ無衝突階層、を
全て導く（出力 S6.1--S6.5）。本章で \texttt{cmbcore} の発展方程式系（P, N, M）が
すべて「導出済みの式」になる。本書の心臓部。

\subsection{Liouville 演算子（導出契約 S6.1）}
光子の位相空間分布 $f(x^\mu,P^\mu)$ は無衝突なら測地線に沿って一定
（$df/d\lambda=0$）。これを変数 $(\eta,x^i,p,\hat p^i)$ で書き下す。質量殻
$P^\mu P_\mu=0$ と計量\eqref{eq:metric}（共形時間版 $ds^2=a^2[-(1+2\Psi)d\eta^2+
(1+2\Phi)dx^2]$）から、固有運動量の大きさ $p$（局所慣性系で測ったエネルギー
$E=pc$）と方向 $\hat p^i$ を定義すると、$P^0=\frac{p}{a}(1-\Psi)$、
$P^i=\frac{p}{a}\hat p^i(1-\Phi)$。全微分は
\begin{equation}
  \frac{df}{d\eta}=\frac{\partial f}{\partial\eta}
  +\frac{dx^i}{d\eta}\frac{\partial f}{\partial x^i}
  +\frac{dp}{d\eta}\frac{\partial f}{\partial p}
  +\frac{d\hat p^i}{d\eta}\frac{\partial f}{\partial\hat p^i},
\end{equation}
で、$d\hat p^i/d\eta$ の項は1次で $\partial f/\partial\hat p^i$（背景等方）に掛かるので
2次となり落ちる。測地線方程式 $dP^\mu/d\lambda+\Gamma^\mu_{\alpha\beta}P^\alpha P^\beta=0$
の $\mu=0$ 成分を1次まで展開すると、運動量の変化率
\begin{equation}
  \frac{1}{p}\frac{dp}{d\eta}=-\mathcal H-\Phi'_{,\eta}-ik\mu\,\Psi
  \label{eq:dpdeta}
\end{equation}
を得る（$-\mathcal H$ は宇宙膨張による赤方偏移、$-\Phi'_{,\eta}$ はポテンシャル時間変化、
$-ik\mu\Psi$ は重力ポテンシャル勾配＝重力赤方偏移；Fourier 空間 $\partial_i\to ik_i$、
$\mu=\hat k\cdot\hat p$）。背景 $-\mathcal H$ は黒体性を保ったまま温度を下げる
（第\ref{ch:ch02_thermal}章）ので、ゆらぎは残り2項で駆動される。

\subsection{輝度（温度ゆらぎ）関数と無衝突方程式}
分布を背景プランク分布 $\bar f$ のまわりで温度ゆらぎ $\Theta$ として
$f=\bar f\!\big(p/[1+\Theta(\eta,\mathbf x,\hat p)]\big)$ とパラメータ化する。
$\bar f$ は $p/T$ のみの関数なので、$\Theta=\delta T/T$ は\emph{$p$ に依らない}方向場と
なり、$p\,\partial_p\bar f=-T\,\partial_T\bar f$ を使うと $\delta f=-p\,\partial_p\bar f\,\Theta$。
これを Liouville 方程式に代入し $p$ 依存を消すと、無衝突部分は
\begin{equation}
  \Theta'_{,\eta}+ick\mu\,\Theta=-\Phi'_{,\eta}-ick\mu\,\Psi,
  \label{eq:collisionless}
\end{equation}
（$' _{,\eta}=\partial/\partial\eta$）。両辺を $\mathcal H$ で割り $x=\ln a$ 形にすれば、
左辺第2項が階層間結合 $\frac{ck}{\mathcal H}\mu\Theta$、右辺が重力源 $-\Phi'-\frac{ck}{\mathcal H}\mu\Psi$ となる。

\subsection{多重極展開と階層（導出契約 S6.3 前半）}
$\mu$ の冪は Legendre 漸化 $\mu P_\ell=\frac{\ell+1}{2\ell+1}P_{\ell+1}+\frac{\ell}{2\ell+1}P_{\ell-1}$
で隣接 $\ell$ を結ぶ。$\Theta=\sum_\ell(-i)^\ell(2\ell+1)P_\ell(\mu)\Theta_\ell$ を
\eqref{eq:collisionless} に代入し $P_\ell$ の係数を比較すると、$\mu\Theta$ の項が
$\Theta_{\ell\pm1}$ への結合を生み、衝突項（次小節）を加えて $x=\ln a$ 形で
（`03` §3.1, P1--P4）:
\begin{align}
  \Theta_0' &= -\frac{ck}{\mathcal H}\Theta_1-\Phi', \label{eq:P1}\\
  \Theta_1' &= \frac{ck}{3\mathcal H}\Theta_0-\frac{2ck}{3\mathcal H}\Theta_2
  +\frac{ck}{3\mathcal H}\Psi+\tau'\Big[\Theta_1+\tfrac13 v_b\Big], \label{eq:P2}\\
  \Theta_\ell' &= \frac{\ell\,ck}{(2\ell+1)\mathcal H}\Theta_{\ell-1}
  -\frac{(\ell+1)ck}{(2\ell+1)\mathcal H}\Theta_{\ell+1}
  +\tau'\Big[\Theta_\ell-\tfrac{1}{10}\Pi\,\delta_{\ell2}\Big],\quad(2\le\ell<\ell_{\max}), \label{eq:P3}
\end{align}
ここで $\Theta_0$ は温度ゆらぎ（密度）、$\Theta_1$ は双極子（速度 $v_\gamma=-3\Theta_1$）、
$\Theta_2$ は四重極（粘性）。四重極の散乱源は\textbf{偏光込み}の
$\Pi\equiv\Theta_2+\Theta^P_0+\Theta^P_2$（次小節）。階層は無限に続くので、最高次は
打ち切り（Callin）
\begin{equation}
  \Theta_{\ell_{\max}}'=\frac{ck}{\mathcal H}\Theta_{\ell_{\max}-1}
  -\frac{c\,(\ell_{\max}+1)}{\mathcal H\,\eta}\,\Theta_{\ell_{\max}}
  +\tau'\Theta_{\ell_{\max}},
  \label{eq:P4}
\end{equation}
で閉じる（$\eta$ は共形時間、$\ell_{\max}=8$；$-c(\ell+1)/(\mathcal H\eta)$ 項は
自由ストリーミングの外向き流出を近似する）。

\subsection{トムソン衝突項と偏光（導出契約 S6.3）}
電子静止系での Thomson 散乱は単極へ等方化させ、四重極依存を残す。実験室系への
ブースト1次で Doppler 項 $\mu v_b$ が出る。多重極展開すると衝突項は
\begin{equation}
  C[\Theta]=-\tau'\Big[\Theta_0-\Theta+\mu v_b-\tfrac{1}{2}P_2(\mu)\,\Pi\Big],
\end{equation}
（$\tau'=-cn_e\sigma_T/H<0$）。Thomson 散乱の偏光依存性により四重極源は温度四重極
だけでなく\textbf{E モード偏光} $\Theta^P_\ell$ も含み、$\Pi=\Theta_2+\Theta^P_0+\Theta^P_2$。
偏光は独立の階層
\begin{align}
  \Theta^{P\prime}_0 &= -\frac{ck}{\mathcal H}\Theta^P_1
    +\tau'\Big[\Theta^P_0-\tfrac12\Pi\Big], \label{eq:Pol0}\\
  \Theta^{P\prime}_\ell &= \frac{\ell\,ck}{(2\ell+1)\mathcal H}\Theta^P_{\ell-1}
    -\frac{(\ell+1)ck}{(2\ell+1)\mathcal H}\Theta^P_{\ell+1}
    +\tau'\Big[\Theta^P_\ell-\tfrac{1}{10}\Pi\,\delta_{\ell2}\Big], \label{eq:Poll}
\end{align}
に従う（打ち切りは\eqref{eq:P4} と同形、$\Theta\to\Theta^P$）。$\Pi$ が温度の四重極源に
入ることで\textbf{減衰尾が正しく}なる：偏光込みの Silk 係数 $16/15$ が実現し、偏光無視
（$\Pi\to\Theta_2$）の $8/9$ より強い減衰となって CLASS と一致する（第\ref{ch:ch09_damping}章・
付録E）。\eqref{eq:P1}--\eqref{eq:Poll} は \texttt{cmbcore.perturbations.\_rhs\_full} に文字単位で
実装されている。

\subsection{CDM とバリオン（導出契約 S6.4, S6.5）}
CDM は衝突なしの非相対論流体で2本（M1, M2）。バリオンは電子・陽子をクーロン結合した
単一流体とし、Thomson 散乱による光子との運動量交換項を持つ（M3, M4）:
\begin{align}
  \delta_c' &= \frac{ck}{\mathcal H}v_c-3\Phi', & v_c' &= -v_c-\frac{ck}{\mathcal H}\Psi, \\
  \delta_b' &= \frac{ck}{\mathcal H}v_b-3\Phi', & v_b' &= -v_b-\frac{ck}{\mathcal H}\Psi
  +\tau'R\,(3\Theta_1+v_b),
\end{align}
ここで結合係数 $R=\frac{4\Omega_{\gamma0}}{3\Omega_{b0}}e^{-x}$。運動量保存（作用反作用）
から、バリオンが光子から受ける力は P2 の衝突項 $\tau'[\Theta_1+v_b/3]$ と符号が逆で、
密度比 $\rho_\gamma/\rho_b\propto R$ 倍に増幅されて $\tau'R(3\Theta_1+v_b)$ となる。

\begin{intuition}
\textbf{なぜ CDM・バリオンは2本で閉じ、光子は閉じないのか}。非相対論的物質は圧力が
無視でき、分布は密度（$\ell=0$）と速度（$\ell=1$）の2モーメントで尽きる——高次
モーメント（圧力・粘性）が $\sim(v/c)^2$ で抑制されるからである。一方、光子は相対論的で
圧力 $p=\rho/3$ を持ち、自由ストリーミングが各 $\ell$ から $\ell\pm1$ へ次々とパワーを
運ぶ（式\eqref{eq:P3} の漸化）ため階層が閉じず、無限次まで（実際には $\ell_{\max}$ で
打ち切って）追う必要がある。これが「音響振動（流体）」と「自由ストリーミング（減衰）」を
分ける物理である。
\end{intuition}

\begin{numreality}
\textbf{記号の罠}: この結合係数 $R=4\Omega_{\gamma0}/(3\Omega_{b0})e^{-x}$ は、
第\ref{ch:ch08_acoustic}章の音速で使うバリオン光子比 $R_s=3\rho_b/4\rho_\gamma$ の
\emph{逆数}である。\texttt{cmbcore} では \texttt{R\_coupling} と \texttt{Rs\_baryon\_photon}
で区別する（`03` §2.3）。
\end{numreality}

\subsection{ニュートリノ}
質量ゼロニュートリノは光子と同形だが\emph{衝突項がない}（$\tau'\to0$）。階層 N1--N4 は
P1--P4 で $\Theta\to\mathcal N$、衝突項を落とした形で、$\ell_{\max}^\nu=10$ まで追う。
自由ストリーミングのみで進むため、地平線進入後は速やかに高次へパワーが流れ
（位相混合で）減衰する——これが放射優勢期のポテンシャル減衰（第\ref{ch:ch08_acoustic}章の
radiation driving）に寄与する。質量ありニュートリノは共動エネルギー
$\varepsilon=\sqrt{(qc)^2+(amc^2)^2}$ が運動量 $q$ に依存するため、各 $q$ ごとに階層を
持つ（$q$ グリッド）。本書は\textbf{背景}を厳密に実装し（第\ref{ch:ch13_params}章,
\texttt{massive\_nu.py}）、摂動レベルの $q$ グリッド階層（$P(k)$ 抑制；TT への寄与は小）は
第2版の拡張とする。

\paragraph{まとめ}
$df/d\lambda=C[f]$ から光子階層\eqref{eq:P1}--\eqref{eq:P3}、CDM・バリオン流体、
ニュートリノ階層を導いた。$R$ と $R_s$ の逆数関係を明示した。これらが
\texttt{cmbcore} の状態ベクトルの発展則そのものである。

\paragraph{演習}
(1) P2 の Doppler 項 $\tfrac13 v_b$ をブースト1次から導け。
(2) 運動量保存から M4 の $\tau'R(3\Theta_1+v_b)$ の符号・係数を確認せよ。
(3) NB3 で $\Theta_\ell$ の階層が高次へエネルギーを運ぶ様子を見よ。
